\documentclass[../cours_3e.tex]{subfiles}
\begin{document}
\setcounter{chapter}{13}
\chapter{Thalès et triangles semblables}
\label{ch:thales}
\begin{objectif}
Calculer avec des parallèles, démontrer un parallélisme, associer les côtés de triangles semblables. Prérequis : proportionnalité, angles, aires et triangles égaux.
\end{objectif}
\section{Thalès : une proportion créée par des parallèles}
\begin{theoreme}
Dans un triangle $ABC$, on prend $M\ne A$ sur $(AB)$ et $N\ne A$ sur $(AC)$. Si $(MN)\parallel(BC)$, alors
\[\frac{AM}{AB}=\frac{AN}{AC}=\frac{MN}{BC}.\]
Les rapports portent ici sur des longueurs positives.
\end{theoreme}
\begin{center}
\begin{tikzpicture}[scale=.8]
\coordinate(A)at(0,3.5);\coordinate(B)at(-2,0);\coordinate(C)at(3,0);
\coordinate(M)at(-1,1.75);\coordinate(N)at(1.5,1.75);
\draw(A)--(B)--(C)--cycle;\draw[accent,thick](M)--(N);
\foreach \p/\pos in {A/above,B/below left,C/below right,M/left,N/right}{\node[\pos]at(\p){$\p$};}
\begin{scope}[xshift=9cm]
\coordinate(A2)at(0,0);\coordinate(B2)at(-2,-1.5);\coordinate(C2)at(2,-1.5);
\coordinate(M2)at(1.2,.9);\coordinate(N2)at(-1.2,.9);
\draw(M2)--(B2)--(C2)--(N2);\draw[accent,thick](M2)--(N2);
\node[right]at(A2){$A$};\node[below]at(B2){$B$};\node[below]at(C2){$C$};
\node[above]at(M2){$M$};\node[above]at(N2){$N$};
\end{scope}
\end{tikzpicture}
\end{center}
À gauche, les triangles sont emboîtés ; à droite, la configuration est dite « papillon ». Toujours faire correspondre $A$ à $A$, $M$ à $B$ et $N$ à $C$.
\section{Démonstration par les aires}
Commençons par $M\in[AB]$ et $N\in[AC]$, avec les points distincts des extrémités.
\begin{demonstration}
Les triangles $BMN$ et $CMN$ ont la même base $MN$ et la même hauteur, car $B,C$ sont sur une parallèle à $(MN)$. Ils ont donc la même aire.
En ajoutant l'aire de $AMN$, on obtient
\[\aire(ABN)=\aire(ACM).\]
Les triangles $ABN$ et $ABC$ ont des bases $AN$ et $AC$ sur la même droite et la même hauteur issue de $B$, donc
\[\frac{\aire(ABN)}{\aire(ABC)}=\frac{AN}{AC}.\]
De même $\aire(ACM)/\aire(ABC)=AM/AB$. Ainsi $AM/AB=AN/AC$.

Pour obtenir le troisième rapport, tracer par $N$ la parallèle à $(AB)$ ; elle coupe $[BC]$ en $D$. Le quadrilatère $BMND$ est un parallélogramme, donc $BD=MN$. Appliquer le rapport déjà démontré dans le triangle $CAB$ donne
\[\frac{CD}{CB}=\frac{CN}{CA}.\]
En soustrayant à $1$ :
\[\frac{BD}{BC}=\frac{AN}{AC},\quad\text{donc}\quad
\frac{MN}{BC}=\frac{AN}{AC}.\]
\end{demonstration}
Si $M,N$ sont au-delà de $B,C$, on applique le cas précédent au grand triangle $AMN$ puis on inverse les rapports. Pour le papillon, la symétrie centrale de centre $A$ envoie $M,N$ de l'autre côté sans changer les longueurs et ramène à des triangles emboîtés (éventuellement en échangeant grand et petit). Les cas d'égalité aux extrémités sont immédiats.
\section{Calculer sans inverser les correspondances}
\begin{methode}
Repérer les deux droites sécantes et leur sommet ; vérifier le parallélisme ; écrire les trois rapports dans le même ordre ; choisir les deux rapports contenant la longueur inconnue et trois longueurs connues.
\end{methode}
\begin{exemple}
Dans la configuration emboîtée, $AB=9$, $AC=12$, $BC=15$ et $AM=6$ (en centimètres). Le rapport vaut $AM/AB=2/3$. Donc $AN=8\unite{cm}$ et $MN=10\unite{cm}$.
Si l'on demande $MB$, il faut encore calculer $AB-AM=3\unite{cm}$.
\end{exemple}
\begin{exemple}[Papillon]
$AB=8$, $AC=6$, $BC=10$ et $AM=12$. Alors le rapport vaut $12/8=1{,}5$, donc $AN=9$ et $MN=15$. Ici $A$ est entre $B$ et $M$ : $BM=BA+AM=20$, et non $4$.
\end{exemple}
\section{Réciproque de Thalès}
\begin{theoreme}
Supposons $A,B,M$ alignés et $A,C,N$ alignés sur deux droites distinctes, avec tous les rapports définis. Supposons aussi que $M$ est sur la demi-droite $[AB)$ exactement quand $N$ est sur $[AC)$ : les deux côtés suivent la même configuration. Si
\[\frac{AM}{AB}=\frac{AN}{AC},\]
alors $(MN)\parallel(BC)$.
\end{theoreme}
\begin{demonstration}
Tracer par $M$ la parallèle à $(BC)$ et appeler $N'$ son intersection avec $(AC)$. Le théorème direct donne $AN'/AC=AM/AB=AN/AC$, donc $AN'=AN$. La condition de configuration place $N'$ et $N$ sur la même demi-droite d'origine $A$. Ils sont donc confondus, et $(MN)$ est la parallèle construite.
\end{demonstration}
\begin{attention}
Les rapports de longueurs ne portent pas de signe. Sans la condition de position, l'égalité des rapports pourrait placer un point du mauvais côté de $A$ et ne garantirait pas le parallélisme.
\end{attention}
\begin{exemple}
$AM=4$, $AB=10$, $AN=6$, $AC=15$, dans une configuration compatible. Les deux rapports valent $2/5$ : les droites sont parallèles. Si $AN=7$, les rapports diffèrent : par contraposée du théorème direct, elles ne sont pas parallèles.
\end{exemple}
\section{Triangles semblables}
\begin{definition}
Deux triangles sont \textbf{semblables} si leurs angles sont égaux deux à deux. Deux angles correspondants égaux suffisent, car la somme des trois angles vaut $180\degre$.
\end{definition}
\begin{theoreme}[Caractérisation par les côtés]
Deux triangles sont semblables si et seulement si leurs côtés correspondants sont proportionnels.
\end{theoreme}
\begin{demonstration}
Si les angles sont égaux, déplacer et éventuellement retourner un triangle pour faire coïncider un sommet et les deux demi-droites qui en partent. L'égalité des angles correspondants impose que les côtés opposés soient parallèles. Thalès donne alors la proportionnalité des trois côtés.

Réciproquement, supposons les trois côtés proportionnels, de rapport $k>0$. Dans l'un des angles du premier triangle, construire un triangle de rapport $k$ avec un côté parallèle au côté opposé. Thalès lui donne les trois longueurs du second triangle. Par CCC, il est égal au second triangle ; or ses angles sont ceux du premier, grâce aux parallèles. Les triangles initiaux sont donc semblables.
\end{demonstration}
\begin{exemple}
Si $\widehat A=40\degre$, $\widehat B=65\degre$, alors $\widehat C=75\degre$. Dans un autre triangle, les angles $40\degre$ et $75\degre$ suffisent à établir la similitude. Il faut associer les côtés \emph{opposés aux mêmes angles}, même si les dessins sont tournés.
\end{exemple}
\begin{propriete}
Pour des triangles semblables de rapport de longueurs $k>0$, les périmètres sont dans le rapport $k$ et les aires dans le rapport $k^2$.
\end{propriete}
\begin{demonstration}
Chaque côté est multiplié par $k$, donc leur somme l'est aussi. Les hauteurs sont multipliées par $k$ (elles sont des longueurs correspondantes de la même similitude). Avec $\aire=bh/2$, l'aire est multipliée par $k\times k=k^2$.
\end{demonstration}
\utilite{Mesurer une hauteur par une ombre, utiliser une maquette, calculer des longueurs dans des figures tournées et justifier les rapports trigonométriques. Une égalité d'angles est parfois plus facile à reconnaître qu'un parallélisme direct.}
\begin{retenir}
Thalès utilise des alignements et des parallèles. Sa réciproque utilise des alignements, une condition de position et une égalité de rapports. Les triangles semblables généralisent la proportionnalité à des figures qui ne sont pas forcément emboîtées.
\end{retenir}
\end{document}
